
\section{ĐỀ SỐ \thesection } %star7
\setcounter{bt}{0}
%%%%%%%%%%%%%%% BÀI 1
\begin{bt} Thực hiện phép tính (tính hợp lí nếu có thể)
	\begin{listEX}[2]
		\item  $A=2025\cdot2026-2025^2$;
        \item $B=3^{2025}:9^{1012}-3\cdot3^2$;
        \item $C=(-2025)-23 +2025+73$;
        \item $D=( -900): \{400+[ (-7)^2+ 1]\} +2025^0$.
	\end{listEX}

	
\loigiai{a) $A=2025$ \hfill b) $B=-24$ \hfill c) $C=50$ \hfill d) $D=-1$
}
\end{bt}

%%%%%%%%%%%%%%% BÀI 2
\begin{bt} Tìm sso nguyên $x$
	\begin{listEX}[2]
		\item  $ x+135=280$;
		\item $3.|2x+3|=201$;
        \item $ 15-5\cdot[(x-30)+6^2]=-20$;
        \item $4^{x+1}+2^2=1028.
$
	\end{listEX}
	\loigiai{
            a) $x=145$;\hspace{1 cm}  
            b) $x=-35$ hay $x=32$; \hspace{1 cm}   
            c) $x=1$;  \hspace{1 cm}   
            d)$x=4$;   \hspace{1 cm}    }
\end{bt}



%%%%%%%%%%%%%%% BÀI 3
\begin{bt}%[Đề bài tự luận]
	Số học sinh khối 6 của trường Ngôi Sao khi xếp hàng $10$ em thì thừa $8$ em,  
	xếp hàng $12$ em thì thừa $10$ em,  
	khi xếp hàng $15$ em thì thừa $13$ em nhưng khi xếp hàng $17$ em thì vừa đủ.  
	
	Tính số học sinh khối 6 của trường, biết rằng số học sinh là một số tự nhiên có $3$ chữ số và nhỏ hơn $250$.
	\loigiai{$60$ ngày}
\end{bt}

%%%%%%%%%%%%%%% BÀI 4
\begin{bt} Hai bạn Lan và Điệp cùng học chung một trường THCS nhưng ở hai lớp khác nhau. An cứ 10 ngày lại trực nhật, Điệp cứ 12 ngày lại trực nhật. Lần đầu tiên hai bạn trực nhật vào cùng một ngày. Hỏi sau ít nhất bao nhiêu ngày hai bạn lại cùng trực nhật? 
 \loigiai{
It nhất sau 60 ngày hai bạn lại cùng trực nhật vào 1 ngày. 
}

\end{bt}
%%%%%%%%%%%%%%% BÀI 5
\begin{bt} 
Một mảnh vườn hình chữ nhật (như hình bên) có chiều dài 9m, chiều rộng 7m. Người ta chia mảnh vườn thành bốn khu gồm hai khu hình vuông cạnh 3m, hai khu hình chữ nhật có chiều dài 5m, chiều rộng 3m và chừa lại phần lối đi như hình bên.
\begin{center}
	\begin{tikzpicture}[scale=0.7,line width=0.8pt,line join=round,>=stealth]
		\draw (0,0)coordinate[label=above:$A$](A)--(3,0)coordinate[label=above:$E$](E)--(4,0)coordinate[label=above:$K$](K)--(9,0)coordinate[label=above:$B$](B)--(9,-3)coordinate[label=right:$Q$](Q)--(9,-4)coordinate[label=right:$T$](T)--(9,-7)coordinate[label=below right:$C$](C)--(6,-7)coordinate[label=below:$R$](R)--(5,-7)coordinate[label=below:$J$](J)--(0,-7)coordinate[label=below left:$D$](D)--(0,-4)coordinate[label=left:$H$](H)--(0,-3)coordinate[label=left:$G$](G)--(A);
		\draw (G)--(3,-3)coordinate[label=below:$F$](F)--(E) (K)--(4,-3)coordinate[label=below:$L$](L)--(Q) (T)--(6,-4)coordinate[label=above:$S$](S)--(R) (J)--(5,-4)coordinate[label=above:$I$](I)--(H);
		\draw ($(A)!.5!(E)$) node[above]{$3$ m}
		($(K)!.5!(B)$) node[above]{$5$ m}
		($(B)!.5!(Q)$) node[right]{$3$ m}
		($(T)!.5!(C)$) node[right]{$3$ m}
		($(R)!.5!(C)$) node[below]{$3$ m}
		($(D)!.5!(J)$) node[below]{$5$ m};
	\end{tikzpicture}
\end{center}
 \begin{listEX}
     \item Tính diện tích phần lối đi.
     \item Người ta muốn lát gạch toàn bộ lối đi bằng những viên gạch hình vuông có cạnh bằng 50cm, biết mỗi viên gạch có giá 25 000 đồng. Tính số tiền cần mua gạch để lát hết lối đi đó.
     \item Trên mảnh vườn, người ta trồng tắc thì thu được 25 000 đồng/$m^2$, trồng quýt thì được 30 000 đồng/$m^2$, trồng cam thì được 40 000 đồng/$m^2$. Em hãy chọn phương có doanh thu tốt hơn trong 2 phương án sau:
\\\textbf{Phương án 1:} Trồng quýt trên cả mảnh đất (trừ lối đi).
\\\textbf{Phương án 2:} Trồng tắc trên mảnh đất hình chữ nhật và trồng cam trên mảnh đất hình vuông. 
\end{listEX}

 
 \loigiai{
 a) $15$ m$^2$ \hfill b) $1\,500\,000$ đồng \hfill c) Phương án 2.}\end{bt}
 %%%%%%%%%%%%%%% BÀI 6
\begin{bt} 
Cho biểu đồ cột cho biết thông tin về kết quả học tập của học sinh khối 6 của  trường Star. Từ biểu đồ em hãy cho biết:
\begin{listEX}
\item Số học sinh xếp loại học lực giỏi là bao nhiêu học sinh? Số học sinh xếp loại học lực giỏi ít hơn số học sinh xếp loại học lực khá bao nhiêu học sinh?
\item Khối 6 có tất cả bao nhiêu học sinh? 
\end{listEX}
     \begin{center}
          	\begin{tikzpicture}[scale=0.7,font=\small,line join=round,>=stealth]
         		% Kích thước khung vẽ
         		\def\hoanh{10.5}
         		\def\tung{6.8}
         		\def\cotwidth{0.7} % độ rộng cột
         		
         		% Tiêu đề
         		\node[above] at (\hoanh/2,\tung+1.2){\normalsize \bfseries Xếp loại học lực học sinh khối 6};
         	 ;
         		
         		% Trục
         		\draw[->] (0,0) node[left]{0} -- (\hoanh,0);
         		\draw[->] (0,0) -- (0,\tung+0.8);
         		
         		% Vạch và lưới ngang (20..160)
         		\foreach \y in {20,40,...,160}{
         			\draw (-0.12,\y/25) node[left] {$\y$} -- ++(0:0.24);
         			\draw[dashed,gray] (0,\y/25) -- (\hoanh-0.2,\y/25);
         		}
         		
         		% === Các cột (x, giá trị) ===
         		% Dùng \x làm TÂM của mỗi cột
         		\foreach \x/\n in {1/38,3.5/140,6/52,8.5/13}{
         			% Cột tô sọc
         			\fill[pattern=north east lines, pattern color=blue!70] 
         			(\x-\cotwidth/2,0) rectangle ++(\cotwidth,\n/25);
         			% Viền cột
         			\draw[black] (\x-\cotwidth/2,0) rectangle ++(\cotwidth,\n/25);
         			% Số trên cột (giữa)
         			\node[above,black] at (\x,\n/25) {\n};
         		}
         		
         		% === Nhãn trục hoành (giữa cột) ===
         		\foreach \x/\p in {1/{Giỏi},3.5/{Khá},6/{Trung bình},8.5/{Yếu}}{
         			\node[align=center] at (\x,-0.55) {\footnotesize \p};
         		}
         		
         		% === Chú thích (Legend) ===
         		\def\lx{7.7}
         		\def\ly{\tung+0.35}
         		\fill[pattern=north east lines, pattern color=blue!70]
         		(\lx-0.35,\ly-0.15) rectangle ++(0.7,0.7);
         		\draw[black] (\lx-0.35,\ly-0.15) rectangle ++(0.7,0.7);
         		\node[anchor=west] at (\lx+0.6,\ly+0.2) {\footnotesize Số học sinh};
         	\end{tikzpicture}
         
         
         
         
     \end{center}
 \loigiai{a) $38$ học sinh, $102$ học sinh \hfill b) $243$ học sinh }
 \end{bt}
  %%%%%%%%%%%%%%% BÀI 7
    \begin{bt} \
    \begin{listEX}
    \item Tìm số tự nhiên n biết rằng $4n+9$ chia hết cho $n-3$ .
    \item Chứng minh $1+4^1+4^2+4^3+...+4^{2023}+4^{2024}$ chia hết cho 21 .
    \end{listEX}
 \loigiai{ 
 \begin{listEX}
     \item $n\in \{4,6,10,24\}$ 
	\item 	Đặt $A=1+4^1+4^2+4^3+...+4^{2023}+4^{2024}$\\
 			$A=(1+4^1+4^2)+(4^3+4^4+4^5)+...+(4^{2022}+4^{2023}+4^{2024})$ \\ 
			$A=(1+4^1+4^2)+4^3.(1+4^1+4^2)+...+4^{2022}(1+4^1+4^2)$\\
			$A= 21+4^3.21+...+4^{2022}.21$\\
			Vậy $A \vdots 21$ 
\end{listEX} 
}
 \end{bt}